\section{Future work}
\label{sec:future_work}

Three concrete strands, plus a longer roadmap of explicitly deferred
work.

\subsection{Concrete next strands}
\label{subsec:future_concrete}

The \textsf{G10} cross-section resolver landed during the most
recent iteration --- the LexScript-borrowed Tarjan-SCC analysis
(\S\ref{subsec:scc_eval}), the lam4-style \texttt{is\_infringed}
predicate, and the Catala-style \texttt{apply\_scope} composition
primitive are all shipped --- so the remaining work in this strand
is the elaboration of the evaluator-side semantics of
\texttt{apply\_scope} into a full element-graph executor that returns
bindings, not just a structural ``does this scope exist'' predicate.

\paragraph{Cross-jurisdiction generalisation.} 8 representative
Indian Penal Code 1860 sections (s300 / s302 / s375 / s376 / s378 /
s420 / s497 / s509) are encoded as a proof-of-concept demonstrating
grammar fit, with a fully-scraped IPC corpus (493 sections) and a
\texttt{yuho refs --compare-libraries} structural-overlap tool
shipped alongside. Full IPC encoding (493 sections) and a richer
SCC-overlap / amendment-divergence analysis are deferred to the
agent-dispatch pipeline. The structural argument and the
falsification criteria are laid out in \S\ref{sec:applicability};
the Akoma Ntoso transpiler is the on-ramp for any cross-jurisdiction
work that wants to interoperate with LegalDocML-consuming tooling.

\paragraph{Deeper precedent integration.} Lifting case-law holding
strings into structured constraints that modify element burden
qualifiers when the matching predicate fires. Right now the L3
stamp records caselaw anchors as opaque strings; the next step is
making them load-bearing.

\paragraph{Evaluator-side \texttt{apply\_scope}.} Elaborating the
cross-section composition primitive into a full element-graph
executor that returns bindings, not just a structural ``does this
scope exist'' predicate.

\subsection{Roadmap (Phase 2A / 2B / 2C)}
\label{subsec:future_roadmap}

The Yuho corpus and tooling extend along three orthogonal axes that
we treat as deferred but explicit, not unscoped.

\paragraph{Phase 2A --- diachronic depth.} The Penal Code 1871 has
150+ years of amendments. Phase 2A extends every encoded section
with every historical version since 1872; \texttt{yuho check
--as-of 1995-06-15 my\_facts.yh} becomes a first-class temporal
query. Acceptance: 524 sections $\times$ at-least-3 historical
versions encoded; per-section timeline rendering on
\texttt{/s/<N>.html}; novel paper contributions on encoding
lineage, an amendment-friction Levenshtein metric, and stable
kernels classifying sections by drift since 1872.

\paragraph{Phase 2B --- vision pivot.} Yuho moves beyond the Penal
Code into the wider Singapore criminal-justice corpus: Evidence
Act 1893, Criminal Procedure Code, Constitution 1963, and Contract
Law (common-law plus Misrepresentation Act and the Contracts
(Rights of Third Parties) Act). Civil remedies and equitable
doctrines need AST shapes Yuho does not yet have, so this is
gated on grammar extension first.

\paragraph{Phase 2C --- interactive learn-by-doing.} The static
explorer site grows into a guided learning surface for law
students: per-section JS fact-builders, in-browser evaluators,
worked-example walkthroughs from the canonical illustrations.
Co-authoring with a law-school educator who can run a classroom
trial unlocks measured learning-outcome evidence; the pedagogical
claim is then ``interactive structural pedagogy for common-law
penal codes''.

\medskip

Phases 2A--2C are deliberately out of scope for the present
artefact: 2A is a multi-month diachronic re-encode, 2B is a
multi-statute scope expansion, and 2C is a co-authored pedagogical
study. The current paper's empirical claims rest entirely on the
static, current-version, single-Act corpus; the roadmap is
provided for honest scoping rather than as load-bearing future
work.
